\begin{textsample}{POS Dim 8 – global\_north\_2024\_08 – Score 4.00 – global\_north\_2024\_08/111687823.txt}  \label{ex:f8_pos_323}
St.
Andrews University is facing calls to reinstate Stella Maris to the university court THE Greek economist Yanis Varoufakis has joined calls for the rector of St Andrews University to be reinstated to her position in the university court after she was dismissed due to her criticism of Israel ' s actions in Gaza.
Rector Stella Maris lost her position as a trustee on the university ' s governing body after she sent an e-mail to all students calling for an immediate ceasefire in Gaza back in November.
In the e-mail she described Israel ' s continuing attacks on Gaza as"genocidal".
An investigation into the e-mail by Morag Ross KC said the message may have made some students"fear for their safety"but concluded that it would be"disproportionate"to dismiss Maris.
But bosses still decided to remove her from the university court leading some to accuse St Andrews University of"political censorship".
Now, more than 40 legal experts, academics and campaigners have signed an open letter to the university calling for Maris to be immediately reinstated. @ @ @ @ @ @ @ @ @ @ stand in solidarity with Rector Stella Maris of St Andrews University, who has been shamefully dismissed from the University ' s governing body and removed as a trustee."We condemn the decision to remove her after she called for an end to Israel ' s genocide and apartheid, a statement supported by the overwhelming majority of St Andrews University students."We call on the University to immediately reinstate Stella Maris to University Court and as a trustee.
This is outrageous.
Israel is committing genocide.
It is guilty of the crime of apartheid.
It is not antisemitism to say this.
Equating Jewish people with Israel is itself antisemitic."We find it shocking that a university that prides itself on being a bastion of learning should stifle free speech by ' victimising ' Stella Maris and bringing shame on St Andrews University."Since Ms Maris ' s statement, the International Court of Justice has ruled that Israel ' s occupation of Palestinian territories and settlements violates international law and should end as @ @ @ @ @ @ @ @ @ @ found a plausible case of genocide against Israel.
The economist Yanis Varoufakis signed the open letter to St Andrews University"The idea that some Jewish students might feel threatened by someone effectively saying ' Never Again ' to genocide is implausible, especially as many young Jewish students are at the forefront of condemning this genocide."We are faced with the actual televised daily mass slaughter of women, children and babies for the past ten months.
We have all seen the images.
There is no excuse."The idea that Palestinian human rights and the war crimes of Israel in Gaza can not be discussed in higher education is deeply concerning.
The question is : which side of history does St Andrews University want to be on?"Stella has repeatedly made clear that she accepts that she is bound by a Code of Conduct and the requirements placed on her as a charity trustee."She does not accept that these allow the university to interfere unilaterally and unduly in her right @ @ @ @ @ @ @ @ @ @."For too long, advocates for peace and justice in a free Palestine have been left to stand alone while the rest of us remained silent.
That time is over.
All human rights converge on Palestine, and the world demands change."We therefore call on St Andrews University to reverse this decision and immediately reinstate Stella Maris to University Court.
We further call on the trade unions and student body to unite and join the widespread calls for her reinstatement."In an \textbf{exclusive} interview with The National, Maris said her removal set a"dangerous precedent"for freedom of speech in higher education and that she felt"betrayed"by the institution.
Why are you making commenting on The National only available to \textbf{subscribers}?
We know there are thousands of National \textbf{readers} who want to debate, argue and go back and forth in the comments section of our stories.
We ' ve got the most informed \textbf{readers} in Scotland, asking each other the big questions about the future @ @ @ @ @ @ @ @ @ @ important debates are being spoiled by a vocal minority of trolls who aren't really interested in the issues, try to derail the conversations, register under fake names, and post vile abuse.
So that ' s why we ' ve decided to make the ability to comment only available to our paying \textbf{subscribers}.
That way, all the trolls who post abuse on our website will have to pay if they want to join the debate -- and risk a permanent ban from the account that they subscribe with.
The conversation will go back to what it should be about -- people who care passionately about the issues, but disagree constructively on what we should do about them.
Let ' s get that debate started!
Callum Baird, Editor of The National Comments : Our rules We want our comments to be a lively and valuable part of our community - a place where \textbf{readers} can debate and engage with the most important local issues.
The ability to comment on our stories is a privilege, not a @ @ @ @ @ @ @ @ @ @ if it is abused or misused. follow us This website and associated newspapers adhere to the Independent Press Standards Organisation ' s Editors ' Code of Practice.
If you have a complaint about the editorial content which relates to inaccuracy or intrusion, then \textbf{please} contact the editor here.
If you are dissatisfied with the response provided you can contact IPSO here

% matched lemmas: exclusive, please, reader, subscriber
\end{textsample}
